% Estilo común de E01 y E02: la letra (Inter, SIL OFL 1.1, en fuentes/) y los colores por rol del tablero
% (generado/colores.tex, que sale de dash_app/tokens.py). Ningún color se escribe aquí: solo roles.
\usepackage{fontspec}
% Sin las alternativas contextuales de Inter (calt): sus glifos de guion, dos puntos y paréntesis «de mayúscula» no
% tienen Unicode, y el texto copiado o extraído del PDF los perdía (2025-11-01 salía 20251101).
\defaultfontfeatures{Ligatures=TeX, RawFeature={-calt}}
\setmainfont{Inter-Regular.ttf}[Path=fuentes/, BoldFont=Inter-SemiBold.ttf, ItalicFont=Inter-Regular.ttf,
  ItalicFeatures={FakeSlant=0.16}, BoldItalicFont=Inter-SemiBold.ttf, BoldItalicFeatures={FakeSlant=0.16}]
\setsansfont{Inter-Regular.ttf}[Path=fuentes/, BoldFont=Inter-SemiBold.ttf, ItalicFont=Inter-Regular.ttf,
  ItalicFeatures={FakeSlant=0.16}]
\newfontfamily\interligera{Inter-Regular.ttf}[Path=fuentes/]
\newfontfamily\intermedia{Inter-Medium.ttf}[Path=fuentes/]
\usepackage[spanish,es-noshorthands,es-tabla,es-nodecimaldot]{babel}
\usepackage[table]{xcolor}
% Generado por documentos/datos_documentos.py desde dash_app/tokens.py (roles de DISENO.md §2).
\definecolor{rol-real}{RGB}{20,20,20}
\definecolor{rol-hoy}{RGB}{133,133,133}
\definecolor{rol-contexto}{RGB}{200,200,200}
\definecolor{rol-modelo}{RGB}{36,87,240}
\definecolor{rol-modelo-final}{RGB}{36,87,240}
\definecolor{rol-previsto-modelo}{RGB}{36,87,240}
\definecolor{rol-previsto-hoy}{RGB}{133,133,133}
\definecolor{rol-recorta}{RGB}{184,35,95}
\definecolor{rol-oraculo}{RGB}{20,20,20}
\definecolor{rol-modelo-tinte}{RGB}{220,229,253}
\definecolor{rol-hoy-tinte}{RGB}{227,227,227}
\definecolor{rol-recorta-tinte}{RGB}{246,219,230}
\definecolor{rol-referencia}{RGB}{102,102,102}
\definecolor{rol-cero}{RGB}{102,102,102}
\definecolor{rol-rejilla}{RGB}{235,235,235}
\definecolor{rol-eje}{RGB}{208,208,208}
\definecolor{rol-silencio}{RGB}{13,148,136}
\definecolor{rol-cra54}{RGB}{232,89,12}
\definecolor{rol-cl72}{RGB}{112,72,232}
\definecolor{rol-bien}{RGB}{24,117,58}
\definecolor{rol-atencion}{RGB}{178,111,11}
\definecolor{rol-atencion-texto}{RGB}{146,91,8}
\definecolor{rol-critico}{RGB}{186,43,46}
\definecolor{rol-acento}{RGB}{36,87,240}
\definecolor{rol-acento-fuerte}{RGB}{27,68,200}
\definecolor{rol-acento-suave}{RGB}{228,235,254}
\definecolor{rol-foco}{RGB}{36,87,240}
\definecolor{rol-tinta}{RGB}{20,20,20}
\definecolor{rol-tinta-2}{RGB}{85,85,85}
\definecolor{rol-tinta-3}{RGB}{102,102,102}
\definecolor{rol-tinta-tenue}{RGB}{138,138,138}
\definecolor{rol-tinta-inversa}{RGB}{255,255,255}
\definecolor{rol-blanco}{RGB}{255,255,255}
\definecolor{rol-superficie}{RGB}{250,250,250}
\definecolor{rol-superficie-2}{RGB}{233,233,233}
\definecolor{rol-fondo}{RGB}{239,239,239}
\definecolor{rol-linea}{RGB}{227,227,227}
\definecolor{rol-linea-fuerte}{RGB}{208,208,208}
\definecolor{rol-pantalla-inicio}{RGB}{20,20,20}
\definecolor{rol-pantalla-pronostico}{RGB}{20,20,20}
\definecolor{rol-pantalla-horario}{RGB}{20,20,20}
\definecolor{rol-pantalla-modelo}{RGB}{20,20,20}
\definecolor{rol-pantalla-resultado}{RGB}{20,20,20}
\definecolor{rol-pantalla-enlaces}{RGB}{20,20,20}

% Generado por documentos/datos_documentos.py: no editar a mano. Cada cifra con su evidencia.

\newcommand{\ModeloActivo}{hgb\_fv2\_20250831\_s42}% api_frecuencias/docs/evidencias/modelo/evaluacion_modelo.json
\newcommand{\PruebaDesde}{2025-11-01}% api_frecuencias/docs/evidencias/modelo/evaluacion_modelo.json
\newcommand{\PruebaHasta}{2025-12-31}% api_frecuencias/docs/evidencias/modelo/evaluacion_modelo.json
\newcommand{\Nprueba}{3.111}% api_frecuencias/docs/evidencias/modelo/evaluacion_modelo.json
\newcommand{\WAPEact}{14,77}% api_frecuencias/docs/evidencias/modelo/evaluacion_modelo.json
\newcommand{\MAEact}{96,27}% api_frecuencias/docs/evidencias/modelo/evaluacion_modelo.json
\newcommand{\WAPEactIC}{13,10--16,87}% api_frecuencias/docs/evidencias/modelo/evaluacion_modelo.json
\newcommand{\WAPEfvuno}{15,07}% api_frecuencias/docs/evidencias/modelo/evaluacion_modelo.json
\newcommand{\MAEfvuno}{98,20}% api_frecuencias/docs/evidencias/modelo/evaluacion_modelo.json
\newcommand{\WAPEfvunoIC}{13,22--17,43}% api_frecuencias/docs/evidencias/modelo/evaluacion_modelo.json
\newcommand{\WAPEbuno}{21,18}% api_frecuencias/docs/evidencias/modelo/evaluacion_modelo.json
\newcommand{\MAEbuno}{138,06}% api_frecuencias/docs/evidencias/modelo/evaluacion_modelo.json
\newcommand{\WAPEbunoIC}{16,61--26,79}% api_frecuencias/docs/evidencias/modelo/evaluacion_modelo.json
\newcommand{\WAPEbdos}{26,40}% api_frecuencias/docs/evidencias/modelo/evaluacion_modelo.json
\newcommand{\MAEbdos}{172,05}% api_frecuencias/docs/evidencias/modelo/evaluacion_modelo.json
\newcommand{\WAPEbdosIC}{22,77--30,57}% api_frecuencias/docs/evidencias/modelo/evaluacion_modelo.json
\newcommand{\SesgoAct}{2,75}% api_frecuencias/docs/evidencias/modelo/evaluacion_modelo.json
\newcommand{\DifActBunoIC}{−10,43 a −3,21}% api_frecuencias/docs/evidencias/modelo/evaluacion_modelo.json
\newcommand{\PmejoraActBuno}{1,00}% api_frecuencias/docs/evidencias/modelo/evaluacion_modelo.json
\newcommand{\MejoraRelBuno}{30}% api_frecuencias/docs/evidencias/modelo/evaluacion_modelo.json
\newcommand{\CobModelo}{79,49}% api_frecuencias/docs/evidencias/modelo/evaluacion_modelo.json
\newcommand{\CobConformal}{77,88}% api_frecuencias/docs/evidencias/modelo/evaluacion_modelo.json
\newcommand{\BajoPdiez}{13,28}% api_frecuencias/docs/evidencias/modelo/evaluacion_modelo.json
\newcommand{\SobrePnoventa}{8,84}% api_frecuencias/docs/evidencias/modelo/evaluacion_modelo.json
\newcommand{\WAPEPico}{14,38}% api_frecuencias/docs/evidencias/modelo/evaluacion_modelo.json
\newcommand{\WAPEValle}{14,97}% api_frecuencias/docs/evidencias/modelo/evaluacion_modelo.json
\newcommand{\WAPERutaSilencio}{15,33}% api_frecuencias/docs/evidencias/modelo/evaluacion_modelo.json
\newcommand{\WAPERutaSilencioBuno}{21,30}% api_frecuencias/docs/evidencias/modelo/evaluacion_modelo.json
\newcommand{\WAPERutaCra}{13,04}% api_frecuencias/docs/evidencias/modelo/evaluacion_modelo.json
\newcommand{\WAPERutaCraBuno}{19,79}% api_frecuencias/docs/evidencias/modelo/evaluacion_modelo.json
\newcommand{\WAPERutaCl}{19,23}% api_frecuencias/docs/evidencias/modelo/evaluacion_modelo.json
\newcommand{\WAPERutaClBuno}{25,27}% api_frecuencias/docs/evidencias/modelo/evaluacion_modelo.json
\newcommand{\RobDosmilveinticuatro}{16,26}% api_frecuencias/docs/evidencias/modelo/evaluacion_modelo.json
\newcommand{\RobDosmilveinticuatroIC}{14,97--17,91}% api_frecuencias/docs/evidencias/modelo/evaluacion_modelo.json
\newcommand{\RobDosmilveinticuatroBuno}{21,55}% api_frecuencias/docs/evidencias/modelo/evaluacion_modelo.json
\newcommand{\RobDosmilveinticinco}{14,48}% api_frecuencias/docs/evidencias/modelo/evaluacion_modelo.json
\newcommand{\RobDosmilveinticincoIC}{13,78--15,28}% api_frecuencias/docs/evidencias/modelo/evaluacion_modelo.json
\newcommand{\RobDosmilveinticincoBuno}{20,04}% api_frecuencias/docs/evidencias/modelo/evaluacion_modelo.json
\newcommand{\CobVdos}{77,82}% api_frecuencias/docs/evidencias/modelo/calibracion_condicional.json
\newcommand{\BajoPdiezVdos}{12,95}% api_frecuencias/docs/evidencias/modelo/calibracion_condicional.json
\newcommand{\SobrePnoventaVdos}{9,23}% api_frecuencias/docs/evidencias/modelo/calibracion_condicional.json
\newcommand{\AnchoVdos}{303,4}% api_frecuencias/docs/evidencias/modelo/calibracion_condicional.json
\newcommand{\CobVuno}{77,88}% api_frecuencias/docs/evidencias/modelo/calibracion_condicional.json
\newcommand{\BajoPdiezVuno}{13,28}% api_frecuencias/docs/evidencias/modelo/calibracion_condicional.json
\newcommand{\SobrePnoventaVuno}{8,84}% api_frecuencias/docs/evidencias/modelo/calibracion_condicional.json
\newcommand{\AnchoVuno}{312,9}% api_frecuencias/docs/evidencias/modelo/calibracion_condicional.json
\newcommand{\CalibPeorDesv}{4,15}% api_frecuencias/docs/evidencias/modelo/calibracion_condicional.json
\newcommand{\HorTreinta}{19,43}% api_frecuencias/docs/evidencias/modelo/horizontes.json
\newcommand{\HorTreintaBuno}{23,12}% api_frecuencias/docs/evidencias/modelo/horizontes.json
\newcommand{\HorSesenta}{14,71}% api_frecuencias/docs/evidencias/modelo/horizontes.json
\newcommand{\HorSesentaBuno}{18,96}% api_frecuencias/docs/evidencias/modelo/horizontes.json
\newcommand{\HorCientoveinte}{10,87}% api_frecuencias/docs/evidencias/modelo/horizontes.json
\newcommand{\HorCientoveinteBuno}{15,44}% api_frecuencias/docs/evidencias/modelo/horizontes.json
\newcommand{\HorEvaluado}{2025-03-01 a 2025-12-31}% api_frecuencias/docs/evidencias/modelo/horizontes.json
\newcommand{\AlertasAhora}{0,53}% api_frecuencias/docs/evidencias/alertas/backtest_alertas.json
\newcommand{\AlertasAntes}{6,08}% api_frecuencias/docs/evidencias/alertas/backtest_alertas.json
\newcommand{\AlertasObsTodas}{1,24}% api_frecuencias/docs/evidencias/alertas/backtest_alertas.json
\newcommand{\DiasRuta}{324}% api_frecuencias/docs/evidencias/alertas/backtest_alertas.json
\newcommand{\AlertasDesde}{2025-09-15}% api_frecuencias/docs/evidencias/alertas/backtest_alertas.json
\newcommand{\AlertasHasta}{2025-12-31}% api_frecuencias/docs/evidencias/alertas/backtest_alertas.json
\newcommand{\LiftCinco}{4,69}% api_frecuencias/docs/evidencias/alertas/backtest_alertas.json
\newcommand{\LiftAntesCinco}{1,77}% api_frecuencias/docs/evidencias/alertas/backtest_alertas.json
\newcommand{\PrecCinco}{0,36}% api_frecuencias/docs/evidencias/alertas/backtest_alertas.json
\newcommand{\LiftSeis}{8,30}% api_frecuencias/docs/evidencias/alertas/backtest_alertas.json
\newcommand{\LiftAntesSeis}{1,31}% api_frecuencias/docs/evidencias/alertas/backtest_alertas.json
\newcommand{\PrecSeis}{0,39}% api_frecuencias/docs/evidencias/alertas/backtest_alertas.json
\newcommand{\LiftSiete}{1,34}% api_frecuencias/docs/evidencias/alertas/backtest_alertas.json
\newcommand{\LiftAntesSiete}{0,78}% api_frecuencias/docs/evidencias/alertas/backtest_alertas.json
\newcommand{\PrecSiete}{0,27}% api_frecuencias/docs/evidencias/alertas/backtest_alertas.json
\newcommand{\LiftOcho}{3,78}% api_frecuencias/docs/evidencias/alertas/backtest_alertas.json
\newcommand{\LiftAntesOcho}{2,87}% api_frecuencias/docs/evidencias/alertas/backtest_alertas.json
\newcommand{\PrecOcho}{0,70}% api_frecuencias/docs/evidencias/alertas/backtest_alertas.json
\newcommand{\UOcupCritica}{0,90}% api_frecuencias/docs/evidencias/alertas/backtest_alertas.json
\newcommand{\UBajaUtil}{0,30}% api_frecuencias/docs/evidencias/alertas/backtest_alertas.json
\newcommand{\UCaida}{0,30}% api_frecuencias/docs/evidencias/alertas/backtest_alertas.json
\newcommand{\UPersistencia}{2}% api_frecuencias/docs/evidencias/alertas/backtest_alertas.json
\newcommand{\UErrorPron}{0,15}% api_frecuencias/docs/evidencias/alertas/backtest_alertas.json
\newcommand{\UZ}{3,5}% api_frecuencias/docs/evidencias/alertas/backtest_alertas.json
\newcommand{\UTolDespacho}{5}% api_frecuencias/docs/evidencias/alertas/backtest_alertas.json
\newcommand{\UTolSalidas}{0,50}% api_frecuencias/docs/evidencias/alertas/backtest_alertas.json
\newcommand{\UAtraso}{2}% api_frecuencias/docs/evidencias/alertas/backtest_alertas.json
\newcommand{\UWapeMax}{0,20}% api_frecuencias/docs/evidencias/alertas/backtest_alertas.json
\newcommand{\USemanasRef}{8}% api_frecuencias/docs/evidencias/alertas/backtest_alertas.json
\newcommand{\UTolSalidasPct}{50}% api_frecuencias/docs/evidencias/alertas/backtest_alertas.json
\newcommand{\UUnoMenosCaida}{0,70}% api_frecuencias/docs/evidencias/alertas/backtest_alertas.json
\newcommand{\AlertasAhoraDosmil}{0,45}% api_frecuencias/docs/evidencias/alertas/backtest_alertas_corte_2024-12-31.json
\newcommand{\AlertasAntesDosmil}{5,70}% api_frecuencias/docs/evidencias/alertas/backtest_alertas_corte_2024-12-31.json
\newcommand{\AlertasPronUno}{0,44}% api_frecuencias/docs/evidencias/alertas/backtest_alertas_pronostico.json
\newcommand{\AlertasPronDos}{0,32}% api_frecuencias/docs/evidencias/alertas/backtest_alertas_pronostico.json
\newcommand{\AlertasPronTres}{0,38}% api_frecuencias/docs/evidencias/alertas/backtest_alertas_pronostico.json
\newcommand{\AlertasPronCuatro}{0,03}% api_frecuencias/docs/evidencias/alertas/backtest_alertas_pronostico.json
\newcommand{\AlertasPronostico}{1,16}% api_frecuencias/docs/evidencias/alertas/backtest_alertas_pronostico.json
\newcommand{\AlertasTotal}{2,40}% api_frecuencias/docs/evidencias/alertas/backtest_alertas_pronostico.json
\newcommand{\ViablesBacktest}{46,0}% api_frecuencias/docs/evidencias/alertas/backtest_alertas_pronostico.json
\newcommand{\RevisionDoble}{dos revisores automatizados con un protocolo escrito, conciliación y kappa de Cohen de 0,80}% api_frecuencias/docs/evidencias/alertas/revision_doble/conciliacion.json
\newcommand{\Kpnoventaycinco}{451}% api_frecuencias/docs/evidencias/pruebas/k6_consultas.json
\newcommand{\Kpnoventaynueve}{572}% api_frecuencias/docs/evidencias/pruebas/k6_consultas.json
\newcommand{\Kpeticiones}{3.553}% api_frecuencias/docs/evidencias/pruebas/k6_consultas.json
\newcommand{\Ktasa}{58,8}% api_frecuencias/docs/evidencias/pruebas/k6_consultas.json
\newcommand{\Kerrores}{0,00}% api_frecuencias/docs/evidencias/pruebas/k6_consultas.json
\newcommand{\Kvus}{50}% api_frecuencias/docs/evidencias/pruebas/k6_consultas.json
\newcommand{\Kingesta}{722}% api_frecuencias/docs/evidencias/pruebas/k6_ingesta.json
\newcommand{\IngestaFilas}{111.792}% api_frecuencias/docs/evidencias/rendimiento.txt
\newcommand{\IngestaSegundos}{4,07}% api_frecuencias/docs/evidencias/rendimiento.txt
\newcommand{\IngestaTasa}{27.469}% api_frecuencias/docs/evidencias/rendimiento.txt
\newcommand{\IngestaMemoria}{53}% api_frecuencias/docs/evidencias/rendimiento.txt
\newcommand{\Ntotal}{1.060}% api_frecuencias/docs/evidencias/pruebas/suite_final.xml
\newcommand{\Npasan}{1.055}% api_frecuencias/docs/evidencias/pruebas/suite_final.xml
\newcommand{\Nfallan}{0}% api_frecuencias/docs/evidencias/pruebas/suite_final.xml
\newcommand{\Nomitidas}{5}% api_frecuencias/docs/evidencias/pruebas/suite_final.xml
\newcommand{\NUnit}{567}% api_frecuencias/docs/evidencias/pruebas/suite_final.xml
\newcommand{\NInteg}{106}% api_frecuencias/docs/evidencias/pruebas/suite_final.xml
\newcommand{\NReplica}{185}% api_frecuencias/docs/evidencias/pruebas/suite_final.xml
\newcommand{\NContrato}{178}% api_frecuencias/docs/evidencias/pruebas/suite_final.xml
\newcommand{\NDatos}{24}% api_frecuencias/docs/evidencias/pruebas/suite_final.xml
\newcommand{\NTcont}{8}% api_frecuencias/docs/evidencias/pruebas/ci/pytest.xml
\newcommand{\NTcontNota}{ (en su trabajo del CI)}% api_frecuencias/docs/evidencias/pruebas/ci/pytest.xml
\newcommand{\FechaPruebas}{2026-11-09 12:37}% api_frecuencias/docs/evidencias/pruebas/suite_final.xml
\newcommand{\CorridaPruebas}{corrida final de la suite completa (2026-11-09 12:37)}% api_frecuencias/docs/evidencias/pruebas/suite_final.xml
\newcommand{\Cobertura}{94,3}% api_frecuencias/docs/evidencias/pruebas/coverage_final.xml
\newcommand{\NIntentosLector}{32}% api_frecuencias/docs/evidencias/pruebas/suite_final.xml
\newcommand{\NIntentosEscritor}{84}% api_frecuencias/docs/evidencias/pruebas/suite_final.xml
\newcommand{\NIntentosTotal}{116}% api_frecuencias/docs/evidencias/pruebas/suite_final.xml
\newcommand{\FuenteIntentos}{\path{api_frecuencias/docs/evidencias/pruebas/suite_final.xml}}% api_frecuencias/docs/evidencias/pruebas/suite_final.xml
\newcommand{\EstadoCI}{en verde en GitHub Actions (run 42627652053, 2026-11-09)}% api_frecuencias/docs/evidencias/pruebas/ci_github.txt
\newcommand{\EstadoCICorto}{en verde en GitHub Actions (run 42627652053, 2026-11-09)}% api_frecuencias/docs/evidencias/pruebas/ci_github.txt
\newcommand{\GateCandidato}{hgb\_fv1\_20250831\_s42}% api_frecuencias/docs/evidencias/monitoreo/compuerta_ejemplo.json
\newcommand{\GateWAPECand}{12,97}% api_frecuencias/docs/evidencias/monitoreo/compuerta_ejemplo.json
\newcommand{\GateWAPEAct}{13,02}% api_frecuencias/docs/evidencias/monitoreo/compuerta_ejemplo.json
\newcommand{\GateDecision}{aprobado}% api_frecuencias/docs/evidencias/monitoreo/compuerta_ejemplo.json
\newcommand{\GateVentana}{2025-09-01 a 2025-09-14}% api_frecuencias/docs/evidencias/monitoreo/compuerta_ejemplo.json
\newcommand{\NVariables}{23}% api_frecuencias/docs/evidencias/monitoreo/onnx_paridad.json
\newcommand{\FeatureVersion}{fv2-2026-11-09}% api_frecuencias/docs/evidencias/monitoreo/onnx_paridad.json
\newcommand{\OnnxDifMax}{0,0028}% api_frecuencias/docs/evidencias/monitoreo/onnx_paridad.json
\newcommand{\OnnxMsFila}{0,10}% api_frecuencias/docs/evidencias/monitoreo/onnx_paridad.json
\newcommand{\SklearnMsFila}{12,6}% api_frecuencias/docs/evidencias/monitoreo/onnx_paridad.json
\newcommand{\FranjasTotal}{13}% api_frecuencias/docs/evidencias/optimizacion/sensibilidad.md
\newcommand{\FranjasViables}{11}% api_frecuencias/docs/evidencias/optimizacion/sensibilidad.md
\newcommand{\FranjasMilpIgual}{13}% api_frecuencias/docs/evidencias/optimizacion/sensibilidad.md
\newcommand{\FranjasNoViables}{CRA 54 (laboral, 7~h) y CRA 54 (sábado, 7~h)}% api_frecuencias/docs/evidencias/optimizacion/sensibilidad.md
\newcommand{\PlanRuta}{CRA 54}% api_frecuencias/docs/evidencias/optimizacion/ejemplo_plan_d5.json
\newcommand{\PlanFecha}{2025-11-19}% api_frecuencias/docs/evidencias/optimizacion/ejemplo_plan_d5.json
\newcommand{\PlanSalidas}{247}% api_frecuencias/docs/evidencias/optimizacion/ejemplo_plan_d5.json
\newcommand{\PlanProgramadas}{290,2}% api_frecuencias/docs/evidencias/optimizacion/ejemplo_plan_d5.json
\newcommand{\PlanHoraPico}{5}% api_frecuencias/docs/evidencias/optimizacion/ejemplo_plan_d5.json
\newcommand{\PlanHeadwayPico}{2,0}% api_frecuencias/docs/evidencias/optimizacion/ejemplo_plan_d5.json
\newcommand{\PlanSalidasPico}{30}% api_frecuencias/docs/evidencias/optimizacion/ejemplo_plan_d5.json
\newcommand{\PlanHeadwayMax}{8,5}% api_frecuencias/docs/evidencias/optimizacion/ejemplo_plan_d5.json
\newcommand{\SonarResumen}{SonarQube: puerta de calidad aprobada (vulnerabilidades: 0; errores: 0; bloqueantes: 0; críticos: 0; duplicación: 0,0~\%; cobertura: 94,3~\%)}% api_frecuencias/docs/evidencias/sonarqube/measures.json
\newcommand{\SonarCorto}{SonarQube: puerta de calidad aprobada (vulnerabilidades: 0; errores: 0; bloqueantes: 0; críticos: 0)}% api_frecuencias/docs/evidencias/sonarqube/measures.json
\newcommand{\SonarPuerta}{puerta de calidad aprobada}% api_frecuencias/docs/evidencias/sonarqube/measures.json
\newcommand{\SonarPruebas}{1.055}% api_frecuencias/docs/evidencias/sonarqube/measures.json
\newcommand{\DemoFilasInicial}{49.623}% api_frecuencias/docs/evidencias/despliegue/demo_ciclo_api.txt
\newcommand{\DemoFilasNuevas}{51}% api_frecuencias/docs/evidencias/despliegue/demo_ciclo_api.txt
\newcommand{\DemoFilasRepeticion}{0}% api_frecuencias/docs/evidencias/despliegue/demo_ciclo_api.txt
\newcommand{\DemoMarcaAntes}{2025-08-30}% api_frecuencias/docs/evidencias/despliegue/demo_ciclo_api.txt
\newcommand{\DemoMarcaDespues}{2025-08-31}% api_frecuencias/docs/evidencias/despliegue/demo_ciclo_api.txt
\newcommand{\DemoWapeNuevo}{26,5}% api_frecuencias/docs/evidencias/despliegue/demo_ciclo_api.txt
\newcommand{\DemoCandidato}{hgb\_fv2\_20250930\_s42}% api_frecuencias/docs/evidencias/despliegue/demo_ciclo_api.txt
\newcommand{\DemoEntrenoSeg}{37}% api_frecuencias/docs/evidencias/despliegue/demo_ciclo_api.txt
\newcommand{\DemoWapeCand}{14,11}% api_frecuencias/docs/evidencias/despliegue/demo_ciclo_api.txt
\newcommand{\DemoWapeAct}{13,98}% api_frecuencias/docs/evidencias/despliegue/demo_ciclo_api.txt
\newcommand{\DemoVentana}{2025-12-01 a 2025-12-14}% api_frecuencias/docs/evidencias/despliegue/demo_ciclo_api.txt
\newcommand{\DemoDecision}{rechazado}% api_frecuencias/docs/evidencias/despliegue/demo_ciclo_api.txt
\newcommand{\DemoTotalSeg}{190}% api_frecuencias/docs/evidencias/despliegue/demo_ciclo_api.txt
\newcommand{\PlanSemanasConexion}{1--4}% PLAN_EVOLUCION.md
\newcommand{\PlanSemanasSombra}{5--8}% PLAN_EVOLUCION.md
\newcommand{\PlanSemanasPiloto}{9--16}% PLAN_EVOLUCION.md
\newcommand{\PlanSemanasSombraTexto}{las semanas 5 a 8}% PLAN_EVOLUCION.md
\newcommand{\PlanSemanasPilotoTexto}{las semanas 9 a 16}% PLAN_EVOLUCION.md
\newcommand{\PlanHallazgosUrgentes}{H3, H6, H7, H9, H13 y H14}% PLAN_EVOLUCION.md
\newcommand{\PlanCriterioSombra}{WAPE \ensuremath{\le} 15 \% a 60 min, \ensuremath{\le} 20 \% a 30 min y \ensuremath{\le} 15 \% a 120 min, y por debajo de B1, en las cuatro semanas; menos del 30 \% de alertas descartadas}% PLAN_EVOLUCION.md
\newcommand{\PlanCostoComputo}{50--80}% PLAN_EVOLUCION.md
\newcommand{\PlanCostoBase}{(A) 0 · (B) 70--150}% PLAN_EVOLUCION.md
\newcommand{\PlanCostoReplica}{(A) 0 · (B) 70--150}% PLAN_EVOLUCION.md
\newcommand{\PlanCostoTotalA}{60--105}% PLAN_EVOLUCION.md
\newcommand{\PlanCostoTotalB}{200--405}% PLAN_EVOLUCION.md
\newcommand{\HallazgosRango}{H1--H16}% api_frecuencias/docs/hallazgos_extreme.md
\newcommand{\NHallazgos}{16}% api_frecuencias/docs/hallazgos_extreme.md
\newcommand{\NHallazgosTexto}{dieciséis}% api_frecuencias/docs/hallazgos_extreme.md
\newcommand{\FuentePruebas}{\path{api_frecuencias/docs/evidencias/pruebas/suite_final.xml}}
\newcommand{\FuenteCobertura}{\path{api_frecuencias/docs/evidencias/pruebas/coverage_final.xml}}

% Generado por documentos/tabla_api.py desde la aplicación FastAPI (app.main). No editar a mano.
\newcommand{\NOperaciones}{62}
\newcommand{\NOperacionesVuno}{60}
\newcommand{\VersionAPI}{1.0.0}
\newcommand{\OrigenTablaAPI}{la aplicación FastAPI (app.main)}

% Generado por documentos/datos_documentos.py desde api_frecuencias/docs/evidencias/respuesta/respuesta_2025-11-19.json. No editar a mano.
\newcommand{\RespuestaDisponible}{1}
\newcommand{\RespuestaFecha}{2025-11-19}
\newcommand{\RespuestaTipoDia}{Laboral}
\newcommand{\RespuestaModelo}{hgb\_fv2\_20250831\_s42}
\newcommand{\RespuestaFuentePronostico}{modelo\ activo}
\newcommand{\RespuestaFuente}{\path{respuesta_2025-11-19.json}}
\newcommand{\RespuestaNoViables}{1}
\newcommand{\RespuestaHoras}{51}
\newcommand{\RespuestaFilas}{%
SILENCIO & 135 / 158,5 & 8,18 / 8,00 / 9,97\newline\textcolor{rol-tinta-3}{hoy 6,86 / 6,77 / 7,26} & ninguna & 11.469 (8.458--14.160) & 44 / 45 & 36,5 / 31,3 \\
CRA 54 & 247 / 290,2 & 3,00 / 4,00 / 7,46\newline\textcolor{rol-tinta-3}{hoy 2,87 / 3,50 / 3,93} & 5~h & 22.800 (19.054--26.560) & 79 / 79 & 72,5 / 52,9 \\
CL 72 & 110 / 96,5 & 10,00 / 9,41 / 9,97\newline\textcolor{rol-tinta-3}{hoy 10,13 / 11,61 / 12,03} & ninguna & 7.525 (5.320--9.628) & 30 / 45 & 23,9 / 23,9 \\
}
\newcommand{\RespEjRuta}{CRA 54}
\newcommand{\RespEjSalidas}{247}
\newcommand{\RespEjProgramadas}{290,2}
\newcommand{\RespEjPicoManana}{3,00}
\newcommand{\RespEjValle}{4,00}
\newcommand{\RespEjPicoTarde}{7,46}
\newcommand{\RespEjNoViables}{a las 5~h la hora sale \emph{no viable}, con su motivo y la mejor aproximación ejecutable}
\newcommand{\RespuestaEjemplo}{CRA 54: ALT-006 crítica, «En 12:00-20:00 (8 h): salieron 17 de 136 despachos esperados (60 / intervalo programado).» Recomendación: Verificar la disponibilidad del vehículo y la operación del despacho.}
\newcommand{\RespuestaEjemploFuente}{\path{respuesta_2025-11-19.json}}

\usepackage{graphicx}
\graphicspath{{figuras/}{../informe/figuras/}}
\usepackage{booktabs,array,tabularx,longtable,multirow}
\usepackage{enumitem}
\usepackage{amsmath}
\usepackage{xurl}
\usepackage[hidelinks]{hyperref}
\usepackage{caption}
\usepackage{tikz}
\usetikzlibrary{positioning,arrows.meta,fit,backgrounds,calc}
\usepackage[most]{tcolorbox}
\usepackage{titlesec}
\usepackage{fancyhdr}
\usepackage{lastpage}

\color{rol-tinta}
\hypersetup{colorlinks=true, linkcolor=rol-acento-fuerte, urlcolor=rol-acento-fuerte, citecolor=rol-acento-fuerte,
  pdfauthor={Abrahan Basto Martínez, Rubén Esguerra Fernández, Paula Gómez}, pdflang={es-CO},
  pdfsubject={Data Challenge 2026 - Extreme Technologies y Universidad del Norte}}
\captionsetup{font={small,color=rol-tinta-2}, labelfont={color=rol-tinta}, justification=raggedright,
  singlelinecheck=false}
\setlist{nosep, leftmargin=1.2em}
\setlength{\parskip}{3pt plus 1pt}
\setlength{\parindent}{0pt}
\renewcommand{\arraystretch}{1.12}
\arrayrulecolor{rol-linea-fuerte}
\widowpenalty=10000 \clubpenalty=10000
\setlength{\emergencystretch}{3em}   % los identificadores largos (\texttt) bajan de línea en vez de salirse
% Texto alineado a la izquierda, como el tablero: los identificadores largos no fuerzan espacios enormes ni se salen.
\usepackage{ragged2e}
\AtBeginDocument{\RaggedRight\hyphenpenalty=5000\exhyphenpenalty=5000}   % casi sin guiones, como el tablero

% Títulos en peso normal (zentra: título grande en 400) y el número en gris.
\titleformat{\section}{\interligera\Large\color{rol-tinta}}{\textcolor{rol-tinta-3}{\thesection}}{0.6em}{}
\titleformat{\subsection}{\intermedia\normalsize\color{rol-tinta}}{\textcolor{rol-tinta-3}{\thesubsection}}{0.5em}{}
\titleformat{\paragraph}[runin]{\intermedia\color{rol-tinta}}{}{0pt}{}[.]
\titlespacing*{\section}{0pt}{10pt plus 2pt}{4pt}
\titlespacing*{\subsection}{0pt}{7pt plus 2pt}{2pt}
\titlespacing*{\paragraph}{0pt}{4pt}{0.6em}

% Fuente de una cifra, tabla o figura: la ruta de la evidencia, relativa a EntregaFinal/.
\newcommand{\fuente}[1]{{\footnotesize\color{rol-tinta-3}Fuente: \path{#1}}}
\newcommand{\evid}[1]{{\footnotesize\color{rol-tinta-3}\path{#1}}}
\newcommand{\evidt}[2]{{\footnotesize\color{rol-tinta-3}\path{#1}, #2}}
\newcommand{\dato}[1]{\textcolor{rol-acento-fuerte}{#1}}
\newcommand{\tenue}[1]{\textcolor{rol-tinta-2}{#1}}

% Tarjeta zentra: superficie casi blanca, borde fino y esquinas redondeadas.
\newtcolorbox{tarjeta}[1][]{enhanced, breakable, colback=rol-superficie, colframe=rol-linea, boxrule=0.6pt, arc=8pt,
  left=8pt, right=8pt, top=6pt, bottom=6pt, #1}
\newtcolorbox{respuesta}[1][]{enhanced, colback=rol-acento-suave, colframe=rol-acento-suave, boxrule=0pt, arc=8pt,
  left=9pt, right=9pt, top=6pt, bottom=6pt, #1}

\pagestyle{fancy}
\fancyhf{}
\renewcommand{\headrulewidth}{0pt}
\fancyfoot[L]{\footnotesize\color{rol-tinta-3}\piedepagina}
\fancyfoot[R]{\footnotesize\color{rol-tinta-3}\thepage\ de \pageref*{LastPage}}
